% Options for packages loaded elsewhere
\PassOptionsToPackage{unicode}{hyperref}
\PassOptionsToPackage{hyphens}{url}
\documentclass[
]{article}
\usepackage{xcolor}
\usepackage{amsmath,amssymb}
\setcounter{secnumdepth}{-\maxdimen} % remove section numbering
\usepackage{iftex}
\ifPDFTeX
  \usepackage[T1]{fontenc}
  \usepackage[utf8]{inputenc}
  \usepackage{textcomp} % provide euro and other symbols
\else % if luatex or xetex
  \usepackage{unicode-math} % this also loads fontspec
  \defaultfontfeatures{Scale=MatchLowercase}
  \defaultfontfeatures[\rmfamily]{Ligatures=TeX,Scale=1}
\fi
\usepackage{lmodern}
\ifPDFTeX\else
  % xetex/luatex font selection
\fi
% Use upquote if available, for straight quotes in verbatim environments
\IfFileExists{upquote.sty}{\usepackage{upquote}}{}
\IfFileExists{microtype.sty}{% use microtype if available
  \usepackage[]{microtype}
  \UseMicrotypeSet[protrusion]{basicmath} % disable protrusion for tt fonts
}{}
\makeatletter
\@ifundefined{KOMAClassName}{% if non-KOMA class
  \IfFileExists{parskip.sty}{%
    \usepackage{parskip}
  }{% else
    \setlength{\parindent}{0pt}
    \setlength{\parskip}{6pt plus 2pt minus 1pt}}
}{% if KOMA class
  \KOMAoptions{parskip=half}}
\makeatother
\usepackage{color}
\usepackage{fancyvrb}
\newcommand{\VerbBar}{|}
\newcommand{\VERB}{\Verb[commandchars=\\\{\}]}
\DefineVerbatimEnvironment{Highlighting}{Verbatim}{commandchars=\\\{\}}
% Add ',fontsize=\small' for more characters per line
\newenvironment{Shaded}{}{}
\newcommand{\AlertTok}[1]{\textcolor[rgb]{1.00,0.00,0.00}{\textbf{#1}}}
\newcommand{\AnnotationTok}[1]{\textcolor[rgb]{0.38,0.63,0.69}{\textbf{\textit{#1}}}}
\newcommand{\AttributeTok}[1]{\textcolor[rgb]{0.49,0.56,0.16}{#1}}
\newcommand{\BaseNTok}[1]{\textcolor[rgb]{0.25,0.63,0.44}{#1}}
\newcommand{\BuiltInTok}[1]{\textcolor[rgb]{0.00,0.50,0.00}{#1}}
\newcommand{\CharTok}[1]{\textcolor[rgb]{0.25,0.44,0.63}{#1}}
\newcommand{\CommentTok}[1]{\textcolor[rgb]{0.38,0.63,0.69}{\textit{#1}}}
\newcommand{\CommentVarTok}[1]{\textcolor[rgb]{0.38,0.63,0.69}{\textbf{\textit{#1}}}}
\newcommand{\ConstantTok}[1]{\textcolor[rgb]{0.53,0.00,0.00}{#1}}
\newcommand{\ControlFlowTok}[1]{\textcolor[rgb]{0.00,0.44,0.13}{\textbf{#1}}}
\newcommand{\DataTypeTok}[1]{\textcolor[rgb]{0.56,0.13,0.00}{#1}}
\newcommand{\DecValTok}[1]{\textcolor[rgb]{0.25,0.63,0.44}{#1}}
\newcommand{\DocumentationTok}[1]{\textcolor[rgb]{0.73,0.13,0.13}{\textit{#1}}}
\newcommand{\ErrorTok}[1]{\textcolor[rgb]{1.00,0.00,0.00}{\textbf{#1}}}
\newcommand{\ExtensionTok}[1]{#1}
\newcommand{\FloatTok}[1]{\textcolor[rgb]{0.25,0.63,0.44}{#1}}
\newcommand{\FunctionTok}[1]{\textcolor[rgb]{0.02,0.16,0.49}{#1}}
\newcommand{\ImportTok}[1]{\textcolor[rgb]{0.00,0.50,0.00}{\textbf{#1}}}
\newcommand{\InformationTok}[1]{\textcolor[rgb]{0.38,0.63,0.69}{\textbf{\textit{#1}}}}
\newcommand{\KeywordTok}[1]{\textcolor[rgb]{0.00,0.44,0.13}{\textbf{#1}}}
\newcommand{\NormalTok}[1]{#1}
\newcommand{\OperatorTok}[1]{\textcolor[rgb]{0.40,0.40,0.40}{#1}}
\newcommand{\OtherTok}[1]{\textcolor[rgb]{0.00,0.44,0.13}{#1}}
\newcommand{\PreprocessorTok}[1]{\textcolor[rgb]{0.74,0.48,0.00}{#1}}
\newcommand{\RegionMarkerTok}[1]{#1}
\newcommand{\SpecialCharTok}[1]{\textcolor[rgb]{0.25,0.44,0.63}{#1}}
\newcommand{\SpecialStringTok}[1]{\textcolor[rgb]{0.73,0.40,0.53}{#1}}
\newcommand{\StringTok}[1]{\textcolor[rgb]{0.25,0.44,0.63}{#1}}
\newcommand{\VariableTok}[1]{\textcolor[rgb]{0.10,0.09,0.49}{#1}}
\newcommand{\VerbatimStringTok}[1]{\textcolor[rgb]{0.25,0.44,0.63}{#1}}
\newcommand{\WarningTok}[1]{\textcolor[rgb]{0.38,0.63,0.69}{\textbf{\textit{#1}}}}
\setlength{\emergencystretch}{3em} % prevent overfull lines
\providecommand{\tightlist}{%
  \setlength{\itemsep}{0pt}\setlength{\parskip}{0pt}}
\usepackage{bookmark}
\IfFileExists{xurl.sty}{\usepackage{xurl}}{} % add URL line breaks if available
\urlstyle{same}
\hypersetup{
  hidelinks,
  pdfcreator={LaTeX via pandoc}}

\author{}
\date{}

\begin{document}

\section{Processing: appunti con videolezioni di
affiancamento.}\label{processing-appunti-con-videolezioni-di-affiancamento}

\begin{Shaded}
\begin{Highlighting}[]
\CommentTok{//[PRIMO FILE]}

\CommentTok{//Se premiamo "play" abbiamo la nostra canvas, di default 100x100px}
\CommentTok{//Usiamo la funzione \textquotesingle{}size()\textquotesingle{} pee definire la nostra dimensione di canvas}
\FunctionTok{size}\OperatorTok{(}\DecValTok{500}\OperatorTok{,}\DecValTok{500}\OperatorTok{);}

\CommentTok{//Adesso non abbiamo un colore di canvas, quindi è opportuno definirlo con \textquotesingle{}background(R, G, B) oppure backgroud(valore in scala di grigio)\textquotesingle{}}
  \CommentTok{//background(0, 150, 255);}
\CommentTok{//Possiamo anche usare in strumenti il Color Selector, e possiam prelevare le coordinate RGB, HSB e Esadecimale}
\FunctionTok{background}\OperatorTok{(}\NormalTok{\#}\ErrorTok{9}\NormalTok{DE9FF}\OperatorTok{);}

\CommentTok{/*Come sappiamo, nell\textquotesingle{}immagine processing il nostro punto di partenza sta in alto a sinistra.}
\CommentTok{Quindi l\textquotesingle{}asse x si sposta in maniera crescente verso destra}
\CommentTok{E l\textquotesingle{}asse y si muove in maniera crescente verso il basso}
\CommentTok{E\textquotesingle{} importante, passare alla funzione come coordinate, prima la \textquotesingle{}x\textquotesingle{} e poi la \textquotesingle{}y\textquotesingle{} */}

\CommentTok{/* Noi useremo spesso la funzione dell\textquotesingle{}ellisse, piuttosto che quella del cerchio.}
\CommentTok{Avendo un centro immaginiamo due diametri, che in realtà sono due assi, e come sempre passiamo pima il valore dell\textquotesingle{}asse orizzontale, e dopo quello verticale}
\CommentTok{  ellipse(x\_c, y\_c, d\_x, d\_y) dove le prime due coordinate sono quelle del punto centrale dell\textquotesingle{}ellissa mentre le altre due sono le diametri orizzontali e verticali */}

\CommentTok{//Adesso disegnamo un cerchio sulla nostra canvas}
\FunctionTok{stroke}\OperatorTok{(}\DecValTok{255}\OperatorTok{,}\DecValTok{0}\OperatorTok{,}\DecValTok{0}\OperatorTok{);} \CommentTok{//I parametri equivalgono allo spazio RGB}
\FunctionTok{ellipse}\OperatorTok{(}\DecValTok{250}\OperatorTok{,}\DecValTok{250}\OperatorTok{,}\DecValTok{50}\OperatorTok{,}\DecValTok{50}\OperatorTok{);}  \CommentTok{//Abbiamo disegnato un cerchio al centro del foglio}
\CommentTok{//Di default disegna il cerchio bianco con un bordo nero. E adesso cambieremo il bordo, ossia il tratto, con questa funzione sotto:}
\CommentTok{// stroke(255,0,0);}
\CommentTok{//Lo inserisco sopra \textquotesingle{}ellipse()\textquotesingle{} perché su Processing si usa un metodo procedurale}

\CommentTok{//Ora facciamo una diagonale orizzontale che attraversa la canvas}
\FunctionTok{line}\OperatorTok{(}\DecValTok{0}\OperatorTok{,} \DecValTok{0}\OperatorTok{,} \DecValTok{500}\OperatorTok{,} \DecValTok{500}\OperatorTok{);}
\CommentTok{//Come sempre i punti sono x1,y1, x2,y2, ovvero dove parte e dove finisce la linea}

\CommentTok{//Ipotizziamo di voler cambiare la dimensione della canvas, ad esempio da 500,500 a 700,700}
\CommentTok{//Poichè le nostre figure non appariranno più al centro, ecco cosa possiamo fare per ovviare al problema}
\CommentTok{// line(0, 0, width, height);}
\CommentTok{// ellipse(width/2,height/2,50,50);}
\CommentTok{// le parole chiave width e height adesso equivalgono alle nostre misure e width/2 e height/2 alla metà, cioè al centro}

\CommentTok{//all\textquotesingle{}esame si potrà consultare la guida di riferimento di Processing, si accede anche tramite Help \textgreater{} Reference}
\CommentTok{//Oppure si seleziona una funzione e con tasto destro si va alla Reference che spiega tutto sulla funzione.}

\CommentTok{//Se vogliamo una figura senza bordom, basta fare:}
\CommentTok{// NoStroke(); //In processing i nomi delle funzioni che sono composti da più parole sono Camel Case}

\CommentTok{//Se invece vogliamo un bordo nero basta fare}
\CommentTok{// stroke(0);}

\CommentTok{//Adesso vediamo come utilizzare i colori con le componenti alpha}

\CommentTok{//Continuo nel secondo file dove svolgeremo un esercizio...}
\end{Highlighting}
\end{Shaded}

\begin{Shaded}
\begin{Highlighting}[]
\CommentTok{//[SECONDO FILE]}

\FunctionTok{size}\OperatorTok{(}\DecValTok{1000}\OperatorTok{,}\DecValTok{700}\OperatorTok{);}
\FunctionTok{background}\OperatorTok{(}\DecValTok{255}\OperatorTok{);}
\FunctionTok{noStroke}\OperatorTok{();}
\CommentTok{//Per cambiare il colore contenuto in una figura, cioè il riempimento si usa \textquotesingle{}fill()\textquotesingle{}}
\FunctionTok{fill}\OperatorTok{(}\DecValTok{255}\OperatorTok{,}\DecValTok{0}\OperatorTok{,}\DecValTok{0}\OperatorTok{,}\DecValTok{160}\OperatorTok{);} \CommentTok{//160 è la trasparenza}
\FunctionTok{ellipse}\OperatorTok{(}\DecValTok{132}\OperatorTok{,}\DecValTok{82}\OperatorTok{,}\DecValTok{200}\OperatorTok{,}\DecValTok{200}\OperatorTok{);}
\FunctionTok{fill}\OperatorTok{(}\DecValTok{0}\OperatorTok{,}\DecValTok{255}\OperatorTok{,}\DecValTok{0}\OperatorTok{,}\DecValTok{160}\OperatorTok{);}
\FunctionTok{ellipse}\OperatorTok{(}\DecValTok{228}\OperatorTok{,{-}}\DecValTok{16}\OperatorTok{,}\DecValTok{200}\OperatorTok{,}\DecValTok{200}\OperatorTok{);} \CommentTok{//Possiamo usare anche coordinate negative}
\FunctionTok{fill}\OperatorTok{(}\DecValTok{0}\OperatorTok{,}\DecValTok{0}\OperatorTok{,}\DecValTok{255}\OperatorTok{,}\DecValTok{160}\OperatorTok{);}    
\FunctionTok{ellipse}\OperatorTok{(}\DecValTok{268}\OperatorTok{,}\DecValTok{118}\OperatorTok{,}\DecValTok{200}\OperatorTok{,}\DecValTok{200}\OperatorTok{);}

\CommentTok{//Oltre alla funzione \textquotesingle{}noStroke()\textquotesingle{} esiste anche la funzione \textquotesingle{}noFill()\textquotesingle{}. Ovviamente inserendo la funzione \textquotesingle{}fill()\textquotesingle{} dopo, invalidiamo ciò che abbiamo scritto.}
\CommentTok{//Quello che farà la funzione noFill() sarà naturalmente quello di NON riempire la figura}

\CommentTok{//Esiste anche la funzione \textquotesingle{}strokeWeight()\textquotesingle{} che di defualt è 1, e il suo scopo è quello di aumentare la dimensione del tratto}

\CommentTok{//Continuo nel terzo file dove svolgeremo un esercizio e presenteremo la funzione rettangolo...}
\end{Highlighting}
\end{Shaded}

\begin{Shaded}
\begin{Highlighting}[]
\CommentTok{//[TERZO FILE]}
\FunctionTok{size}\OperatorTok{(}\DecValTok{1000}\OperatorTok{,} \DecValTok{700}\OperatorTok{);}
\FunctionTok{background}\OperatorTok{(}\DecValTok{255}\OperatorTok{);}
\FunctionTok{noFill}\OperatorTok{();}

\CommentTok{//Funzione per fare rettangoli e quadrati}
  \CommentTok{//rect(x,y,l,h);}
    \CommentTok{//x e y sono del vertice in alto a sinistra}
    \CommentTok{//l e h sono lunghezza e altezza}
\FunctionTok{rect}\OperatorTok{(}\DecValTok{50}\OperatorTok{,}\DecValTok{60}\OperatorTok{,}\DecValTok{80}\OperatorTok{,}\DecValTok{80}\OperatorTok{);}

\CommentTok{//Adesso facciamo un ellisse e un rettangolo con le stesse coordinate}
\FunctionTok{ellipse}\OperatorTok{(}\DecValTok{50}\OperatorTok{,}\DecValTok{60}\OperatorTok{,}\DecValTok{80}\OperatorTok{,}\DecValTok{80}\OperatorTok{);}

\CommentTok{//Ipotizziamo di voler cerchio inscritto nel quadrato.}
\CommentTok{//Ecco che introduciamo le funzioni di modalità.}

\CommentTok{/*DALLA GUIDA PROCESSING:}
\CommentTok{ellipseMode(CORNER) interpreta i primi due parametri di ellipse() come angolo superiore sinistro della forma, mentre il terzo e il quarto parametro sono la sua larghezza e altezza.}
\CommentTok{(Spiegare meglio)*/}
\FunctionTok{ellipseMode}\OperatorTok{(}\NormalTok{CORNER}\OperatorTok{);}
\FunctionTok{rect}\OperatorTok{(}\DecValTok{150}\OperatorTok{,}\DecValTok{60}\OperatorTok{,}\DecValTok{80}\OperatorTok{,}\DecValTok{80}\OperatorTok{);}
\FunctionTok{ellipse}\OperatorTok{(}\DecValTok{150}\OperatorTok{,}\DecValTok{60}\OperatorTok{,}\DecValTok{80}\OperatorTok{,}\DecValTok{80}\OperatorTok{);}

\CommentTok{//Nel caso della modalità \textquotesingle{}RADIUS\textquotesingle{} cambia il modo di interpretare il terzo ed il quarto parametro, prima erano il diametro, adesso il raggio}
\CommentTok{/*DALLA GUIDA PROCESSING:}
\CommentTok{ellipseMode(RADIUS) utilizza i primi due parametri di ellipse() come punto centrale della forma, ma utilizza il terzo e il quarto parametro per specificare metà della larghezza e dell\textquotesingle{}altezza della forma.}
\CommentTok{(Spiegare meglio)*/}
\FunctionTok{ellipseMode}\OperatorTok{(}\NormalTok{RADIUS}\OperatorTok{);}
\FunctionTok{rect}\OperatorTok{(}\DecValTok{250}\OperatorTok{,}\DecValTok{60}\OperatorTok{,}\DecValTok{80}\OperatorTok{,}\DecValTok{80}\OperatorTok{);}
\FunctionTok{ellipse}\OperatorTok{(}\DecValTok{250}\OperatorTok{,}\DecValTok{60}\OperatorTok{,}\DecValTok{80}\OperatorTok{,}\DecValTok{80}\OperatorTok{);}

\CommentTok{//Come si torna alla modalità di deufault? Con \textquotesingle{}CENTER\textquotesingle{}}
\FunctionTok{ellipseMode}\OperatorTok{(}\NormalTok{CENTER}\OperatorTok{);}
\FunctionTok{rect}\OperatorTok{(}\DecValTok{350}\OperatorTok{,}\DecValTok{60}\OperatorTok{,}\DecValTok{80}\OperatorTok{,}\DecValTok{80}\OperatorTok{);}
\FunctionTok{ellipse}\OperatorTok{(}\DecValTok{350}\OperatorTok{,}\DecValTok{60}\OperatorTok{,}\DecValTok{80}\OperatorTok{,}\DecValTok{80}\OperatorTok{);}

\CommentTok{//L\textquotesingle{}ultima modalità è \textquotesingle{}CORNERS\textquotesingle{}. Diamo angolo in alto a sinistra, angolo in alto a destra, e lui disegna in questa porzione}
\CommentTok{/*DALLA GUIDA DI PROCESSING:}
\CommentTok{ellipseMode(CORNERS) interpreta i primi due parametri di ellipse() come posizione di un angolo del riquadro di delimitazione dell\textquotesingle{}ellisse e il terzo e il quarto parametro come posizione dell\textquotesingle{}angolo opposto.}
\CommentTok{(Spiegare meglio)*/}
\FunctionTok{ellipseMode}\OperatorTok{(}\NormalTok{CORNERS}\OperatorTok{);}
\FunctionTok{rect}\OperatorTok{(}\DecValTok{450}\OperatorTok{,}\DecValTok{60}\OperatorTok{,}\DecValTok{80}\OperatorTok{,}\DecValTok{80}\OperatorTok{);}
\FunctionTok{ellipse}\OperatorTok{(}\DecValTok{450}\OperatorTok{,}\DecValTok{60}\OperatorTok{,}\DecValTok{550}\OperatorTok{,}\DecValTok{160}\OperatorTok{);}

\CommentTok{//Link a tutte le modalità dell\textquotesingle{}ellipseMode che spiega bene e con disegni: https://processing.org/reference/ellipseMode\_.html}

\CommentTok{//Esiste anche \textquotesingle{}rectMode()\textquotesingle{} che fanno quasi la stessa cosa ma per i rettangoli}

\CommentTok{//Passiamo al quarto file...}
\end{Highlighting}
\end{Shaded}

\begin{Shaded}
\begin{Highlighting}[]
\CommentTok{//[QUARTO FILE]}

\FunctionTok{size}\OperatorTok{(}\DecValTok{1000}\OperatorTok{,} \DecValTok{700}\OperatorTok{);}
\CommentTok{//Andiamo a vedere come "conservare" un colore, e si fa con la variabile \textquotesingle{}color\textquotesingle{}}
\NormalTok{color c1 }\OperatorTok{=} \FunctionTok{color}\OperatorTok{(}\DecValTok{204}\OperatorTok{,}\DecValTok{154}\OperatorTok{,}\DecValTok{0}\OperatorTok{);} \CommentTok{//Abbiamo chiamato la funzione \textquotesingle{}color()\textquotesingle{} nella variabile \textquotesingle{}color c1\textquotesingle{} perchhé non possiamo passare tre valori alla funzione, dunque se abbiamo un colore in RGB , dunque con le 3 componenit, lo passiamo alla funzione che si occupa di fare il casting e passarlo alla variabile \textquotesingle{}color\textquotesingle{}}
\NormalTok{color c2 }\OperatorTok{=}\NormalTok{ \#FFCC00}\OperatorTok{;}
\FunctionTok{noStroke}\OperatorTok{();}
\FunctionTok{fill}\OperatorTok{(}\NormalTok{c1}\OperatorTok{);}
\FunctionTok{rect}\OperatorTok{(}\DecValTok{0}\OperatorTok{,}\DecValTok{0}\OperatorTok{,}\DecValTok{40}\OperatorTok{,}\NormalTok{height}\OperatorTok{);}
\FunctionTok{fill}\OperatorTok{(}\NormalTok{c2}\OperatorTok{);}
\FunctionTok{rect}\OperatorTok{(}\DecValTok{40}\OperatorTok{,}\DecValTok{0}\OperatorTok{,}\DecValTok{40}\OperatorTok{,}\NormalTok{height}\OperatorTok{);}
\NormalTok{color c3 }\OperatorTok{=} \FunctionTok{get}\OperatorTok{(}\DecValTok{10}\OperatorTok{,}\DecValTok{10}\OperatorTok{);} \CommentTok{//\textquotesingle{}get()\textquotesingle{} serve per prelevare un colore nella coordinata corrispondente}
\FunctionTok{fill}\OperatorTok{(}\NormalTok{c3}\OperatorTok{);}
\FunctionTok{rect}\OperatorTok{(}\DecValTok{80}\OperatorTok{,}\DecValTok{0}\OperatorTok{,}\DecValTok{40}\OperatorTok{,}\NormalTok{height}\OperatorTok{);}
\end{Highlighting}
\end{Shaded}

\begin{Shaded}
\begin{Highlighting}[]
\CommentTok{//[QUINTO FILE]}

\FunctionTok{size}\OperatorTok{(}\DecValTok{1000}\OperatorTok{,} \DecValTok{700}\OperatorTok{);}
\CommentTok{// strokeWeight(24);}

\CommentTok{//Queste linee hanno le punte arrotondate, e si chiamano Cap, se vogliamo modificarle, aggiungiamo prima delle linee delle funzioni}
\CommentTok{/*}
\CommentTok{    line(60,25,130,95);}
\CommentTok{    line(160,25,230,95);}
\CommentTok{    line(260,25,330,95);}
\CommentTok{    line(360,25,430,95);}

\CommentTok{Per la funzione che si usa è \textquotesingle{}strokeCap()\textquotesingle{}}
\CommentTok{*/}

    \CommentTok{//line(60,25,130,95);}
    \CommentTok{//strokeCap(SQUARE);}
    \CommentTok{//line(160,25,230,95);}
    \CommentTok{//strokeCap(PROJECT);}
    \CommentTok{//line(260,25,330,95);}
    \CommentTok{//strokeCap(ROUND);}
    \CommentTok{//line(360,25,430,95);}

\CommentTok{//Qual è la differenza tra SQUARE e PROJECT?}
\CommentTok{/*}
\CommentTok{Quando noi abbiamo una linea che abbiamo definito con strokeWeight, Processing prende la nostra linea, e se il nostro strokeWeight equivale a 24 allora mette un margine di 12 da un lato e un lato di 12 da un altro. Per consegnarci uno spessore finale di 24.}
\CommentTok{{-} Nel caso di SQUARE allora i cap verranno tagliati}
\CommentTok{{-} Nel caso di PROJECT verranno aggiunti 12px alle estremità, ossia ai cap}
\CommentTok{Infatti una linea PROJECT è più lunga di una SQUARE }
\CommentTok{*/}

\CommentTok{//Una funzione simile si può usare con i rettangoli, che si occupa delle diverse smussature degli spigoli/giunture. I joun sono i punti di giunture nello spigolo}
\CommentTok{//E la funzione che si usa si chiama \textquotesingle{}strokeJoin()\textquotesingle{}}
\CommentTok{//strokeWeight(12);}
\CommentTok{//rect(60,25,70,70);}
\CommentTok{//strokeJoin(ROUND); //Modifica i jouin con una forma rotonda }
\CommentTok{//rect(160,25,70,70);}
\CommentTok{//strokeJoin(BEVEL); //Modifica i join con una forma smussata}
\CommentTok{//rect(260,25,70,70);}
\CommentTok{//strokeJoin(MITER); //Riporta alla forma di defualt del rettangolo, con un forma squadrata (classica forma geometrica rettangolare)}
\CommentTok{//rect(360,25,70,70);}

\CommentTok{//Esiste anche la funzione triangle che richiede 6 parametri ((x,y) per ogni punto), l\textquotesingle{}ordine di inserimento dei vertici nella funzione non ha importaza. }
\FunctionTok{fill}\OperatorTok{(}\DecValTok{21}\OperatorTok{,} \DecValTok{64}\OperatorTok{,} \DecValTok{120}\OperatorTok{);}
\FunctionTok{strokeWeight}\OperatorTok{(}\DecValTok{10}\OperatorTok{);}
\FunctionTok{triangle}\OperatorTok{(}\DecValTok{200}\OperatorTok{,}\DecValTok{400}\OperatorTok{,} \DecValTok{200}\OperatorTok{,}\DecValTok{100}\OperatorTok{,} \DecValTok{500}\OperatorTok{,}\DecValTok{400}\OperatorTok{);}
\end{Highlighting}
\end{Shaded}

\begin{Shaded}
\begin{Highlighting}[]
\CommentTok{//[SESTO FILE]}
\CommentTok{//Introduciamo i costrutti loop e if/else}

\FunctionTok{size}\OperatorTok{(}\DecValTok{800}\OperatorTok{,}\DecValTok{500}\OperatorTok{);}
\ControlFlowTok{for} \OperatorTok{(}\DataTypeTok{int}\NormalTok{ i }\OperatorTok{=} \DecValTok{20}\OperatorTok{;}\NormalTok{ i }\OperatorTok{\textless{}=}\DecValTok{400}\OperatorTok{;}\NormalTok{ i}\OperatorTok{+=}\DecValTok{20}\OperatorTok{)\{}
  \FunctionTok{stroke}\OperatorTok{(}\DecValTok{0}\OperatorTok{);} \CommentTok{//Ristabiliamo lo stroke nero}
  \FunctionTok{line}\OperatorTok{(}\NormalTok{i}\OperatorTok{,}\DecValTok{0}\OperatorTok{,}\NormalTok{i}\OperatorTok{+}\NormalTok{i}\OperatorTok{/}\DecValTok{2}\OperatorTok{,}\NormalTok{height}\OperatorTok{/}\DecValTok{2}\OperatorTok{);} \CommentTok{//Linea superiore}
  \FunctionTok{line}\OperatorTok{(}\NormalTok{i}\OperatorTok{,}\NormalTok{height}\OperatorTok{,}\NormalTok{i}\OperatorTok{+}\NormalTok{i}\OperatorTok{/}\DecValTok{2}\OperatorTok{,}\NormalTok{height}\OperatorTok{/}\DecValTok{2}\OperatorTok{);} \CommentTok{//Linea orizzontale}
  \FunctionTok{stroke}\OperatorTok{(}\DecValTok{255}\OperatorTok{,}\DecValTok{0}\OperatorTok{,}\DecValTok{0}\OperatorTok{);}
  \FunctionTok{line}\OperatorTok{(}\DecValTok{20}\OperatorTok{,}\NormalTok{i}\OperatorTok{,}\DecValTok{60}\OperatorTok{,}\NormalTok{i}\OperatorTok{);}
\OperatorTok{\}}
\end{Highlighting}
\end{Shaded}

\begin{Shaded}
\begin{Highlighting}[]
\CommentTok{//[SETTIMO FILE]}

\FunctionTok{size}\OperatorTok{(}\DecValTok{500}\OperatorTok{,}\DecValTok{500}\OperatorTok{);}
\FunctionTok{background}\OperatorTok{(}\DecValTok{0}\OperatorTok{);}
\FunctionTok{noStroke}\OperatorTok{();}
\DataTypeTok{int}\NormalTok{ d }\OperatorTok{=} \DecValTok{40}\OperatorTok{;}

\ControlFlowTok{for}\OperatorTok{(}\DataTypeTok{int}\NormalTok{ x }\OperatorTok{=}\NormalTok{ d}\OperatorTok{/}\DecValTok{2}\OperatorTok{;}\NormalTok{ x }\OperatorTok{\textless{}=}\NormalTok{ height}\OperatorTok{;}\NormalTok{ x}\OperatorTok{+=}\NormalTok{d}\OperatorTok{)\{}
  \ControlFlowTok{for}\OperatorTok{(}\DataTypeTok{int}\NormalTok{ y }\OperatorTok{=}\NormalTok{ d}\OperatorTok{/}\DecValTok{2}\OperatorTok{;}\NormalTok{ y }\OperatorTok{\textless{}=}\NormalTok{ width}\OperatorTok{;}\NormalTok{ y}\OperatorTok{+=}\NormalTok{d}\OperatorTok{)\{}
     \FunctionTok{ellipse}\OperatorTok{(}\NormalTok{x}\OperatorTok{,}\NormalTok{y}\OperatorTok{,}\NormalTok{ d}\OperatorTok{,}\NormalTok{d}\OperatorTok{);}
 \OperatorTok{\}}
\OperatorTok{\}}
\end{Highlighting}
\end{Shaded}

\begin{Shaded}
\begin{Highlighting}[]
\CommentTok{//[OTTAVO FILE]}
\CommentTok{//Introduzione alla modalità interattiva}
\CommentTok{//Usiamo l\textquotesingle{}override (spiegare meglio)}

\CommentTok{//Introduciamo due funzioni di Processing, di tipo \textquotesingle{}void\textquotesingle{} che sono setup e draw}
\DataTypeTok{void} \FunctionTok{setup}\OperatorTok{()\{} \CommentTok{//Qui mettiamo le istruzioni per l\textquotesingle{}avvio del programma. Viene eseguito una e una sola  volta}
  \CommentTok{//\textquotesingle{}print()\textquotesingle{}, per stampare nella console e il cursore rimane lì}
  \CommentTok{//\textquotesingle{}println()\textquotesingle{}, per stampare nella console ed andare a new line}
  \FunctionTok{println}\OperatorTok{(}\StringTok{"sono nel setup"}\OperatorTok{);} \CommentTok{//Stampato solo una volta}
  \CommentTok{//frameRate(2); //per cambiare il frame rate, scriverà 2 "sono nel draw" per 2 volte al secondo.}
   \FunctionTok{frameRate}\OperatorTok{(}\DecValTok{30}\OperatorTok{);}
   \FunctionTok{background}\OperatorTok{(}\DecValTok{0}\OperatorTok{);}
\OperatorTok{\}}

\CommentTok{//void draw()\{//Qui invece implementiamo un ciclo infinito, come se fosse un while == true. Viene eseguito ad 1 frame di default a 60 frame}
  \CommentTok{//println("sono nel draw "+frameCount+" frame sono passati"); //Stampato 60 volte al secondo //chiarire questa sintassi {-}\textgreater{} "+frameCount+"}
  \CommentTok{//frameCount è una variabile che ci fa sapere quanti fotogrammi sono passati dall\textquotesingle{}inizio dell\textquotesingle{}interazione}
  \CommentTok{//if(frameCount \textgreater{} 150)\{ //passano 5 secondi a 30 frame al secondo}
    \CommentTok{//ellipse(width/2, height/2, 50,50); //Disegna un elisse a dopo 5 secondi}
  \CommentTok{//\}}
\CommentTok{//\}}

\CommentTok{//Questo esercizio, dove viene utilizzato frameCount, può essere fatto usando la funzione \textquotesingle{}millis()\textquotesingle{} al posto di frameCount, è una funzione che restituisce i millisecondi passati dall\textquotesingle{}inizio dell\textquotesingle{}esecuzione del programma.}
\CommentTok{//Commenterò l\textquotesingle{}esercizio con frameCount per mostrare millis()}

\DataTypeTok{void} \FunctionTok{draw}\OperatorTok{()\{}\CommentTok{//Qui invece implementiamo un ciclo infinito, come se fosse un while == true. Viene eseguito ad 1 frame di default a 60 frame}
  \FunctionTok{println}\OperatorTok{(}\StringTok{"sono nel draw "}\OperatorTok{+}\NormalTok{frameCount}\OperatorTok{+}\StringTok{" frame sono passati"}\OperatorTok{);} \CommentTok{//Stampato 60 volte al secondo //chiarire questa sintassi {-}{-}\textgreater{} "+frameCount+"}
  \CommentTok{//frameCount è una variabile che ci fa sapere quanti fotogrammi sono passati dall\textquotesingle{}inizio dell\textquotesingle{}interazione}
  \ControlFlowTok{if}\OperatorTok{(}\FunctionTok{millis}\OperatorTok{()} \OperatorTok{\textgreater{}} \DecValTok{5000}\OperatorTok{)\{} \CommentTok{//impostiamo 5000 perché 5 secondi in millisecondi sonoo 5000 (1000 millisecondi in un secondo)}
    \FunctionTok{ellipse}\OperatorTok{(}\NormalTok{width}\OperatorTok{/}\DecValTok{2}\OperatorTok{,}\NormalTok{ height}\OperatorTok{/}\DecValTok{2}\OperatorTok{,} \DecValTok{50}\OperatorTok{,}\DecValTok{50}\OperatorTok{);} \CommentTok{//Disegna un elisse a dopo 5 secondi}
  \OperatorTok{\}}
\OperatorTok{\}}
\end{Highlighting}
\end{Shaded}

\begin{Shaded}
\begin{Highlighting}[]
\CommentTok{//[NONO FILE]}
\CommentTok{//Renderemo interattivo l\textquotesingle{}esercizio sui cerchi di prima}

\CommentTok{//mettiamo la variabile \textquotesingle{}d\textquotesingle{} fuori dalle funzioni così ho uno scope più ampio e quidi è globale}
  \DataTypeTok{int}\NormalTok{ d }\OperatorTok{=} \DecValTok{40}\OperatorTok{;}

\DataTypeTok{void} \FunctionTok{setup}\OperatorTok{()\{}
  \FunctionTok{size}\OperatorTok{(}\DecValTok{500}\OperatorTok{,}\DecValTok{500}\OperatorTok{);}
  \FunctionTok{background}\OperatorTok{(}\DecValTok{0}\OperatorTok{);}
  \FunctionTok{noStroke}\OperatorTok{();}
\OperatorTok{\}}

\DataTypeTok{void} \FunctionTok{draw}\OperatorTok{()\{}

  \ControlFlowTok{for}\OperatorTok{(}\DataTypeTok{int}\NormalTok{ y }\OperatorTok{=}\NormalTok{ d}\OperatorTok{/}\DecValTok{2}\OperatorTok{;}\NormalTok{ y }\OperatorTok{\textless{}=}\NormalTok{ height}\OperatorTok{;}\NormalTok{ y}\OperatorTok{+=}\NormalTok{d}\OperatorTok{)\{}
    \ControlFlowTok{for}\OperatorTok{(}\DataTypeTok{int}\NormalTok{ x }\OperatorTok{=}\NormalTok{ d}\OperatorTok{/}\DecValTok{2}\OperatorTok{;}\NormalTok{ x }\OperatorTok{\textless{}=}\NormalTok{ width}\OperatorTok{;}\NormalTok{ x}\OperatorTok{+=}\NormalTok{d}\OperatorTok{)\{}
      \ControlFlowTok{if}\OperatorTok{(}\FunctionTok{dist}\OperatorTok{(}\NormalTok{x}\OperatorTok{,}\NormalTok{y}\OperatorTok{,}\NormalTok{ mouseX}\OperatorTok{,}\NormalTok{ mouseY}\OperatorTok{)} \OperatorTok{\textless{}}\NormalTok{ d}\OperatorTok{/}\DecValTok{2}\OperatorTok{)\{} \CommentTok{//usiamo una variabile che tramite coordinate (x,y) mantiene le coordinate del cursore del nostro mouse.}
          \CommentTok{//Tenendo conto del fatto che le variabili del nostro mouse sono mouseX e mouseY, quelle del centro del cerchio nella canvas sono xc e yc (dove c sta per centro). Per verificare la condizione verifichiamo le coordinate del mouse con quelle del centro del cerchio. E quindi scriviamo che la distanza di questi due punti è minore di quella del centro. Usiamo la funzione \textquotesingle{}dist()\textquotesingle{} che fa il calcolo della distanza tra le coordinate e il centro).}
          \CommentTok{//Quindi come condizione abbiamo detto "se la distanza tra questi punti è minore del raggio, allora siamo dentro il cerchio e puoi colorarlo"}
        \FunctionTok{fill}\OperatorTok{(}\DecValTok{0}\OperatorTok{,}\DecValTok{0}\OperatorTok{,}\DecValTok{255}\OperatorTok{);}      
      \OperatorTok{\}}\ControlFlowTok{else}\OperatorTok{\{}
        \FunctionTok{fill}\OperatorTok{(}\DecValTok{255}\OperatorTok{);}
      \OperatorTok{\}}
      \FunctionTok{ellipse}\OperatorTok{(}\NormalTok{x}\OperatorTok{,}\NormalTok{y}\OperatorTok{,}\NormalTok{d}\OperatorTok{,}\NormalTok{d}\OperatorTok{);}
    \OperatorTok{\}}
  \OperatorTok{\}}
\OperatorTok{\}}
\end{Highlighting}
\end{Shaded}

\begin{Shaded}
\begin{Highlighting}[]
\CommentTok{//[DECIMO FILE]}
\CommentTok{//Cerchiamo di rendere l\textquotesingle{}esercizio di prima più interattivo, Facciamo in modo che all\textquotesingle{}aumentare delle coordinate dell\textquotesingle{}asse X aumenteremo la componente di rosso, e all\textquotesingle{}aumentare di quelle di Y avremo l\textquotesingle{}aumento della componente di verde.}
\CommentTok{//Vuol dire che se siamo nell\textquotesingle{}angolo in alto a sinistra avremo nero puro, in alto a destra i rosso puro, il basso a sinistra il verde puro e in basso a destra il giallo puro}

\DataTypeTok{int}\NormalTok{ d }\OperatorTok{=} \DecValTok{40}\OperatorTok{;}

\DataTypeTok{void} \FunctionTok{setup}\OperatorTok{()\{}
  \FunctionTok{size}\OperatorTok{(}\DecValTok{500}\OperatorTok{,}\DecValTok{500}\OperatorTok{);}
  \FunctionTok{background}\OperatorTok{(}\DecValTok{0}\OperatorTok{);}
  \FunctionTok{noStroke}\OperatorTok{();}
  \FunctionTok{colorMode}\OperatorTok{(}\NormalTok{RGB}\OperatorTok{,}\NormalTok{ width}\OperatorTok{);} \CommentTok{//Impostiamo il range [0{-}500] al posto di [0{-}255] e usiamo lo spazio colore RGB}
\OperatorTok{\}}

\DataTypeTok{void} \FunctionTok{draw}\OperatorTok{()\{}
  \ControlFlowTok{for}\OperatorTok{(}\DataTypeTok{int}\NormalTok{ y }\OperatorTok{=}\NormalTok{ d}\OperatorTok{/}\DecValTok{2}\OperatorTok{;}\NormalTok{ y }\OperatorTok{\textless{}=}\NormalTok{ height}\OperatorTok{;}\NormalTok{ y}\OperatorTok{+=}\NormalTok{d}\OperatorTok{)\{}
    \ControlFlowTok{for}\OperatorTok{(}\DataTypeTok{int}\NormalTok{ x }\OperatorTok{=}\NormalTok{ d}\OperatorTok{/}\DecValTok{2}\OperatorTok{;}\NormalTok{ x }\OperatorTok{\textless{}=}\NormalTok{ width}\OperatorTok{;}\NormalTok{ x}\OperatorTok{+=}\NormalTok{d}\OperatorTok{)\{}
      \ControlFlowTok{if}\OperatorTok{(}\FunctionTok{dist}\OperatorTok{(}\NormalTok{x}\OperatorTok{,}\NormalTok{y}\OperatorTok{,}\NormalTok{ mouseX}\OperatorTok{,}\NormalTok{ mouseY}\OperatorTok{)} \OperatorTok{\textless{}}\NormalTok{ d}\OperatorTok{/}\DecValTok{2}\OperatorTok{)\{}
        \FunctionTok{fill}\OperatorTok{(}\NormalTok{x}\OperatorTok{,}\NormalTok{y}\OperatorTok{,}\NormalTok{ width}\OperatorTok{);}      
      \OperatorTok{\}}\ControlFlowTok{else}\OperatorTok{\{}
        \FunctionTok{fill}\OperatorTok{(}\NormalTok{x}\OperatorTok{,}\NormalTok{y}\OperatorTok{,}\DecValTok{0}\OperatorTok{);}
      \OperatorTok{\}}
      \FunctionTok{ellipse}\OperatorTok{(}\NormalTok{x}\OperatorTok{,}\NormalTok{y}\OperatorTok{,}\NormalTok{d}\OperatorTok{,}\NormalTok{d}\OperatorTok{);}
    \OperatorTok{\}}
  \OperatorTok{\}}
\OperatorTok{\}}
\end{Highlighting}
\end{Shaded}

\begin{Shaded}
\begin{Highlighting}[]
\CommentTok{//[UNDICESIMO FILE]
}
\CommentTok{//Ipotizziamo di voler disegnare un poliedro, quindi non una delle semplici figura che abbiamo fatto ora
}
\CommentTok{//Per farlo utilizzamo una funzione \textquotesingle{}beginShape()\textquotesingle{}, in cui inserire all\textquotesingle{}interno tutti i vertici \textquotesingle{}vertex()\textquotesingle{} della nostra figura
}
\CommentTok{//Disegnamo su paint e ci segniamo le coordinate della figura, per evitare i calcoli.
}


\FunctionTok{size}\OperatorTok{(}\DecValTok{640}\OperatorTok{,} \DecValTok{280}\OperatorTok{);}



\FunctionTok{beginShape}\OperatorTok{();} \CommentTok{//unisce i punti forniti successivamente
}
\CommentTok{//per cambiare fill e stroke si fa all\textquotesingle{}interno della funzione \textquotesingle{}beginShape()\textquotesingle{}
}
\FunctionTok{fill}\OperatorTok{(}\DecValTok{153}\OperatorTok{,}\DecValTok{176}\OperatorTok{,}\DecValTok{180}\OperatorTok{);} \CommentTok{//colore di riempimento
}
\FunctionTok{vertex}\OperatorTok{(}\DecValTok{180}\OperatorTok{,}\DecValTok{82}\OperatorTok{);}

\FunctionTok{vertex}\OperatorTok{(}\DecValTok{207}\OperatorTok{,}\DecValTok{36}\OperatorTok{);}

\FunctionTok{vertex}\OperatorTok{(}\DecValTok{214}\OperatorTok{,}\DecValTok{63}\OperatorTok{);}

\FunctionTok{vertex}\OperatorTok{(}\DecValTok{407}\OperatorTok{,}\DecValTok{11}\OperatorTok{);}

\FunctionTok{vertex}\OperatorTok{(}\DecValTok{412}\OperatorTok{,}\DecValTok{30}\OperatorTok{);}

\FunctionTok{vertex}\OperatorTok{(}\DecValTok{219}\OperatorTok{,}\DecValTok{82}\OperatorTok{);}

\FunctionTok{vertex}\OperatorTok{(}\DecValTok{226}\OperatorTok{,}\DecValTok{109}\OperatorTok{);}

\CommentTok{//endShape(); //la figura resta aperta
}
\FunctionTok{endShape}\OperatorTok{(}\NormalTok{CLOSE}\OperatorTok{);} \CommentTok{//chiude la figura (unisce l\textquotesingle{}ultimo vertice col primo)
}


\CommentTok{//avremmo anche potuto riscrivere il primo vertice alla fine, come ultimo vertice}
\end{Highlighting}
\end{Shaded}


\end{document}
